\chapter{Related Work}
\label{ch:related}

\section{Automated Document Verification}

The problem of verifying whether a document conforms to a specification is closely related
to form understanding, a well-studied area in document image analysis. Early work by
Doermann and Rosenfeld~\cite{doermann1995document} proposed a structural description
language for forms, where named fields were associated with spatial anchors. Their system
could identify missing or misaligned fields in scanned government forms. The key limitation
of early approaches was the assumption of fixed-position fields, which fails for label types
whose layout varies with the presence or absence of optional zones.

More recent work has moved toward deep learning. Xu et al.~\cite{xu2020layoutlm} proposed
LayoutLM, which jointly encodes text tokens and their 2D spatial positions using a BERT-style
transformer. LayoutLMv3~\cite{huang2022layoutlmv3} further incorporates image patch
embeddings, achieving state-of-the-art results on form understanding benchmarks such as
FUNSD and CORD. While these models are powerful, they require fine-tuning on
domain-specific labelled data and produce probabilistic outputs that are harder to explain
to a compliance officer than a simple ``field X expected value Y, found Z.''

\section{Industrial Label and Packaging Verification}

Commercial machine vision systems from providers such as Cognex, Keyence, and Zebra
Technologies are widely used in manufacturing lines for label inspection tasks such as:
verifying that a barcode decodes correctly, checking that required symbols (CE mark, WEEE
symbol) are present, and confirming that text fields are not blank. These systems typically
operate in milliseconds using hardware cameras and FPGA-accelerated processing.

However, they are designed for end-of-line inspection of printed labels under controlled
lighting, not for pre-press template validation of digital PDF files or scanned images. They
also lack the ability to check content correctness (e.g.\ does the article number on the label
match the GTIN encoded in the DataMatrix barcode?) because they treat the label as an
image to inspect, not a document to parse.

A closer analogue to our work is the pharmaceutical packaging verification literature.
Mack et al.~\cite{mack2017pharmaceutical} described an automated system for verifying
that drug packaging inserts contain all required text fields and regulatory symbols, using
\ac{OCR} and rule-based matching. Their system achieved~95\% precision on a curated
dataset of 400 drug inserts. The key difference from our setting is that pharmaceutical
packaging tends to have higher-contrast text on lighter backgrounds, and the regulatory
requirements are more standardised than IKEA's proprietary label specifications.

\section{OCR Engine Comparison Studies}

Multiple studies have compared \ac{OCR} engine performance on specialised document
types. Hegghammer~\cite{hegghammer2022ocr} benchmarked Tesseract, ABBYY FineReader,
and Amazon Textract on historical printed texts, finding Tesseract competitive in clean-image
conditions but inferior on degraded or handwritten content. Bukhari et al.~\cite{bukhari2012improved}
studied the effect of image preprocessing on Tesseract's performance on document images,
finding that adaptive thresholding and deskewing can improve character accuracy by up to~12
percentage points.

For industrial labels specifically, the rotated text zones (90° rotations are common in IKEA
labels) and small font sizes (as low as~2mm, corresponding to roughly~7pt at 96~DPI)
present particular challenges for general-purpose \ac{OCR} engines trained on natural
documents. Our evaluation in \cref{ch:evaluation} provides evidence that an ensemble
strategy mitigates these shortcomings more effectively than single-engine tuning.

\section{Barcode Quality and Decoding}

ISO/IEC~15416 and ISO/IEC~15415 define print quality standards for linear and 2D barcodes
respectively, specifying metrics such as modulation, symbol contrast, decodability, and
defects. Automated barcode quality grading is a well-established industrial practice.
Russ and Schuchardt~\cite{russ2020barcode} surveyed software graders for DataMatrix
symbols, finding that many commercial graders achieve~99\%+ agreement with certified
verifiers on clean symbols but diverge on marginal-quality symbols.

Our system does not implement full ISO grading (which would require a physical verifier),
but does check that a DataMatrix can be decoded and that its content satisfies the required
AI structure—a coarser but practical check suitable for pre-press validation.

\section{Explainable Validation Systems}

The tension between accuracy and interpretability is well documented in the \ac{ML}
literature~\cite{lipton2018mythos}. For compliance use cases, interpretability is often
non-negotiable: regulatory frameworks such as the EU AI Act require that high-risk AI
systems produce human-understandable explanations for their decisions.

Ribeiro et al.~\cite{ribeiro2016should} introduced LIME, which generates local surrogate
explanations for any black-box model. Lundberg and Lee~\cite{lundberg2017unified} proposed
SHAP, which provides theoretically grounded feature attribution. Both methods are commonly
applied post-hoc to neural classifiers. In contrast, rule-based validators are \emph{inherently}
interpretable: every decision traces to a specific rule evaluation, making post-hoc explanation
methods unnecessary. Our work therefore deliberately chooses rule-based validation over a
neural classifier, accepting a modest accuracy trade-off in exchange for full traceability.

\section{Positioning of This Work}

Table~\ref{tab:related_comparison} positions this thesis relative to the closest related
systems along four dimensions: whether the system handles PDF-native extraction, whether
it validates content (not just presence), whether it identifies the label type automatically,
and whether its decisions are traceable to specific rules.

\begin{table}[h]
\centering
\caption{Comparison of this work with related systems.}
\label{tab:related_comparison}
\begin{tabularx}{\textwidth}{Xccccc}
\toprule
\textbf{System} & \textbf{PDF-native} & \textbf{Content} & \textbf{Auto type} & \textbf{Traceable} & \textbf{Barcode} \\
\midrule
Cognex ViDi~\cite{cognex2023} & \xmark & \xmark & \xmark & \xmark & \checkmark \\
LayoutLMv3~\cite{huang2022layoutlmv3} & \xmark & \checkmark & \checkmark & \xmark & \xmark \\
Mack et al.~\cite{mack2017pharmaceutical} & \xmark & \checkmark & \xmark & \checkmark & \xmark \\
\textbf{This work} & \checkmark & \checkmark & \checkmark & \checkmark & \checkmark \\
\bottomrule
\end{tabularx}
\end{table}

The combination of PDF-native text extraction, automatic label-type detection, content-level
field validation, barcode parsing, and rule-traceable decisions is, to the best of our
knowledge, novel in the context of industrial label \ac{QA}.
